\documentclass[10pt,a4paper]{amsart}
\usepackage[margin=19mm]{geometry}
\usepackage{amsmath,amssymb,mathtools}
\usepackage[scr=rsfs,cal=euler]{mathalfa}
\usepackage{xfrac}
\usepackage[unicode,hidelinks,bookmarks=false]{hyperref}
\newcommand{\ie}{i.e.\ }
\newcommand{\cf}{cf.\ }
\newcommand{\resp}{respectively\ }
\newcommand{\Subsection}{\subsection}
\providecommand{\nobreakdash}{\nobreak\hbox{-}}
\setlength{\emergencystretch}{2em}
\multlinegap0pt
\usepackage{amssymb}
\usepackage{bm}
\usepackage{dsfont}
\usepackage{csquotes}
\usepackage{mathtools}
\usepackage{braket}
\usepackage{tikz}
\usepackage{stackengine}
\usepackage[shortlabels]{enumitem}
\newtheorem{lemma}{Lemma}
\title{Generalised Complex and Spinor Relations}
\author{Thomas C. De Fraja, Vincenzo Emilio Marotta, Richard J. Szabo}
\date{Source: arXiv:2603.10819v2 (CC BY 4.0)}
\begin{document}
\maketitle

\noindent\emph{Source and licence.} Generalised Complex and Spinor Relations. Version arXiv:2603.10819v2 (CC BY 4.0), distributed by arXiv under the Creative Commons Attribution 4.0 International licence, http://creativecommons.org/licenses/by/4.0/, as recorded in the arXiv licence metadata for this exact version and shown on the abstract page https://arxiv.org/abs/2603.10819v2. Full licence notice: \texttt{licenses/arxiv-2603.10819-CC-BY-4.0.txt}.\par\smallskip

\noindent\textit{Editorial scope.} Counted unit: the article's lemma lem:formsdefined (source0.tex lines 2808--2816). The article's complete original proof of it is delivered with it (source0.tex lines 2818--2853, one \texttt{proof} environment of 36 lines). No mathematics is written, repaired or re-derived: the statement and the proof are the article's own lines, in the article's own notation, and no retained line is substituted or re-flowed.  No further source-local context is needed: every cross-reference of every retained line resolves inside the two delivered spans.  Nothing was omitted from any retained span, so the package carries no omission record.  No retained line contains a citation, so the wrapper carries no bibliography and no bibliography entry.  No retained line contains a figure environment, an image-inclusion command or drawing code, so no image is delivered and the package declares no asset dependency.  Editorial formatting only, in addition: an AMS \texttt{amsart} preamble that restates the article's own declarations and macros used by the retained lines, namely the environments \texttt{lemma} on separate counters, together with the standard packages those lines need.  The retained environments print in this document as Lemma 1 in order of appearance, whereas the article prints the same items with the numbering of its own full document.  The complete native source of the article as distributed in the arXiv e-print for this exact version is bundled unchanged as \texttt{sources/196170/source0.tex}.
\begin{lemma} \label{lem:formsdefined}
The forms defined above satisfy
    \begin{enumerate}
        \item \ $\displaystyle \exp(\widehat \alpha_1) \cdot \frac{1}{m_2!}\,\alpha_1^{m_2} = \exp(\alpha_1) \ , $ \\
        \item \ $\exp(\widehat \alpha_{\rm mix}) \cdot ( f_\alpha^1 \wedge \cdots \wedge f_\alpha^{m_3}) = (\ddot b_\alpha^1 + f_\alpha^1) \wedge \cdots \wedge (\ddot b_\alpha^{m_3} + f_\alpha^{m_3}) \ , $ \\[-3mm]
        \item \ $\widehat \mu_1 \cdot(\mu_1 \wedge \overline{\mu}_1) = \overline{\mu}_1 \ , $ \ and \\[-3mm]
        \item \ $\widehat\mu_{\rm mix} \cdot (f_\mu^1 \wedge \cdots \wedge f_\mu^{k_3}) = \exp(\ddot{b}_\mu^1 \wedge f_\mu^1 + \cdots +\ddot{b}_\mu^{k_3} \wedge f_\mu^{k_3}) \ .$
    \end{enumerate}
\end{lemma}

\begin{proof}
    (1) Firstly, for each natural number $j\leqslant m_2$
    \begin{align}
        \widehat \alpha_1^j \cdot \alpha_1^{m_2} = m_2\, j\ \widehat\alpha_1^{j-1} \cdot \alpha_1^{m_2-1} = \cdots = \frac{m_2!\, j!}{(m_2-j)!}\,\alpha_1^{m_2-j} \ .
    \end{align}
    Thus
    \begin{equation}
    \begin{aligned}
        \exp(\widehat \alpha_1) \cdot \frac{1}{m_2!}\, \alpha_1^{m_2} = \sum_{j=0}^{m_2}\,\frac1{j!}\,\widehat\alpha_1^j  \cdot \frac{1}{m_2!}\, \alpha_1^{m_2}=\sum_{j=0}^{m_2}\, \frac{1}{m_2!\,j!}\, \frac{m_2!\, j!}{(m_2-j)!}\, \alpha_1^{m_2-j} = \exp( \alpha_1)\ .
    \end{aligned}
    \end{equation}
    
    \noindent (2) We compute
    \begin{equation}
    \begin{aligned}
        \exp(\widehat \alpha_{\rm mix}) \cdot ( f_\alpha^1 \wedge \cdots \wedge f_\alpha^{m_3}) &= \exp(\ddot b_\alpha^1 \wedge \widehat f_\alpha^{\,1} ) \wedge \cdots \wedge \exp(\ddot b_\alpha^{m_3} \wedge \widehat f_\alpha^{\,m_3} ) \cdot ( f_\alpha^1 \wedge \cdots \wedge f_\alpha^{m_3})\\[4pt]
        &=\big(\exp(\ddot b_\alpha^1 \wedge \widehat f_\alpha^{\,1} ) \cdot f_\alpha^1\,\big) \wedge \cdots \wedge \big(\exp(\ddot b_\alpha^{m_3} \wedge \widehat f_\alpha^{\,m_3} )\cdot f_\alpha^{m_3}\,\big) \\[4pt]
        &=\big((1+\ddot b_\alpha^1 \wedge \widehat f_\alpha^{\,1} ) \cdot f_\alpha^1 \big)\wedge \cdots \wedge \big((1+\ddot b_\alpha^{m_3} \wedge \widehat f_\alpha^{\,m_3} )\cdot f_\alpha^{m_3}\big)\\[4pt]
        &=(\ddot b_\alpha^1 + f_\alpha^1) \wedge \cdots \wedge (\ddot b_\alpha^{m_3} + f_\alpha^{m_3}) \ ,
    \end{aligned}
    \end{equation}
    where we use the fact that $\widehat f_\alpha^{\,i} \cdot \ddot b_\alpha^j = 0$ for all $i,j$ and $\widehat f_\alpha^{\,i} \cdot f_\alpha^j = 0$ when $i\neq j$ to commute the exponentials.

    \noindent (3) Follows immediately.

    \noindent (4) Similarly to (2) we compute
    \begin{equation}
    \begin{aligned}
        \widehat\mu_{\rm mix} \cdot (f_\mu^1 \wedge \cdots \wedge f_\mu^{k_3}) &= (\ddot{b}_\mu  + \widehat{f}_\mu^{\,1}) \wedge \cdots \wedge (\ddot{b}_\mu^{k_{3}} + \widehat{f}_\mu^{\,k_{3}}) \cdot (f_\mu^1 \wedge \cdots \wedge f_\mu^{k_3})\\[4pt]
        &= \big((\ddot{b}_\mu  + \widehat{f}_\mu^{\,1}) \cdot f_\mu^1\big)\wedge \cdots \wedge \big((\ddot{b}_\mu^{k_{3}} + \widehat{f}_\mu^{\,k_{3}}) \cdot f_\mu^{k_3}\big)\\[4pt]
        &= (1 + \ddot{b}_\mu\wedge  f_\mu^1)\wedge \cdots \wedge (1+ \ddot{b}_\mu^{k_{3}}\wedge f_\mu^{k_3})\\[4pt]
        &= \exp(\ddot{b}_\mu^1 \wedge f_\mu^1 + \cdots + \ddot{b}_\mu^{k_3} \wedge f_\mu^{k_3}) \ ,
    \end{aligned}
    \end{equation}
    as required.
\end{proof}

\end{document}
